\begin{center}
\fontsize{16pt}{14pt}\selectfont
\bfseries Лопатин Александр Сергеевич\mdseries

\fontsize{14pt}{14pt}\selectfont
Резюме
\end{center}
{
\fontsize{12pt}{12pt}\selectfont
\begin{flushright}
\item дата рождения: 10 июня 1990
\item мобильный телефон: +7 911 83 66445
%\item e-mail: \includegraphics[scale=0.2]{email}
\item e-mail: sbar.geek ät gmail.com
\item примеры кода: \
\href{https://github.com/sbar?tab=repositories}{https://github.com/sbar}
\end{flushright}
}

\setcounter{section}{0}
\section{Умения}
\begin{itemize}
\item C и \bfseries C++\mdseries
\subitem ООП, темплейты, основные паттерны проектирования,
фреймворки
\bfseries Qt \mdseries и
\bfseries Marmalade \mdseries
(aka \bfseries AirPlay\mdseries),
базовые знания \bfseries OpenGL \mdseries и \bfseries GLSL\mdseries

\subitem кроссплатформенная разработка
(\bfseries GNU/Linux\mdseries, Windows, iOS, \bfseries Android \mdseries и др.)

\subitem использование средств отладки и анализа (gdb, valgrind, cppcheck),
опыт работы с Visual~Studio

\item небольшой опыт работы с \bfseries Java \mdseries (J2ME, Android)

\item скриптинг
\subitem \bfseries Python \mdseries
(сетевые приложения, парсинг, базы данных и другое)
\subitem \bfseries UNIX shell (bash)\mdseries
\subitem небольшой опыт с \bfseries Perl \mdseries и \bfseries Lua\mdseries

\item SQL (базовые знания, разработка реляционных баз данных)
в реализациях SQLite, MySQL и PostgreSQL
\end{itemize}

\section{Опыт работы}
\begin{itemize}
\item {
\fontsize{14pt}{14pt}\selectfont
\bfseries EA (Electronic Arts)\mdseries, Санкт-Петербург
}
\subitem \bfseries с сентября, 2012 ``Программист''\mdseries
\subsubitem Портирование мобильных игровых приложений
\subsubitem Проект \href{http://www.facebook.com/realracing}{Real Racing 3} ---\
портирование на Android

\item {
\fontsize{14pt}{14pt}\selectfont
\bfseries SPL (SmartPhoneLabs)\mdseries, Великий Новгород
}
\subitem \bfseries с января, 2012 ``Младший инженер-программист''\mdseries
\subsubitem Разработка и портирование компьютерных игр для мобильных платформ
(iOS, Android, BlackBerry, J2ME и др.)
\subsubitem Проекты:
\subsubitem\href{https://itunes.apple.com/us/app/panda-blair!/id500995558?mt=8}{Panda Blair!} --- разработка дополнительных игровых фич, портирование\
на Android и BlackBerry
\subsubitem\href{http://java.mob.org/game/phineas\_and\_ferb\_time\_travelers.html}{Phineas and Ferb: Time Travelers} --- портирование на устройства с J2ME

\subitem \bfseries с сентября, 2011 ``Младший сотрудник службы качества''
\mdseries
\subsubitem ручное тестирование J2ME-приложений
\subsubitem автоматизация внутренней деятельности посредством скриптинга на\
UNIX shell и Python
\end{itemize}

\section{Образование}
\begin{itemize}
%\item с 2011 Северо-Западный Государственный Технический Университет
%(Санкт-Петербург), специальность
%``Программное обеспечение вычислительной техники и автоматизированных систем'',
%заочное~обучение
\item 2005---2009 Многопрофильный политехнический колледж НовГУ, специальность
``Программное обеспечение вычислительной техники и автоматизированных систем''
\end{itemize}

\section{Другое}
Есть умение понимать и изменять чужой код.

Владею английским языком на среднем уровне (достаточно для чтения/написания
документации и делового общения); родной язык --- русский.

Являюсь уверенным пользователем и разработчиком Gentoo GNU/Linux.

Есть опыт использования систем контроля версиями (git в основном) и
багтрекеров (Jira, Bugzilla, Launchpad, Mantis и др.)

Быстрообучаемость новым технологиям.
Прошел военную службу в Вооруженных Силах.
Не~курю.
% vim: textwidth=0
